\documentclass[11pt,a4paper]{article}
\usepackage[margin=2.5cm]{geometry}
\usepackage{graphicx}
\usepackage{booktabs}
\usepackage{amsmath}
\usepackage{amssymb}
\usepackage{xcolor}
\usepackage{hyperref}
\usepackage{caption}
\usepackage{subcaption}
\usepackage{float}
\usepackage{parskip}
\usepackage{titlesec}
\usepackage{enumitem}
\usepackage{multirow}
\usepackage{array}
\usepackage{colortbl}

\definecolor{correct}{RGB}{39,174,96}
\definecolor{wrong}{RGB}{192,57,43}
\definecolor{marginal}{RGB}{230,126,34}
\definecolor{sectionblue}{RGB}{41,128,185}

\titleformat{\section}{\large\bfseries\color{sectionblue}}{}{0em}{}[\titlerule]
\titleformat{\subsection}{\normalsize\bfseries}{}{0em}{}

\hypersetup{colorlinks=true, linkcolor=sectionblue, urlcolor=sectionblue}

\title{
    \textbf{Cascade Directionality Detection via MIL Training Dynamics}\\[0.5em]
    \large Session Experiment Report
}
\author{Cytokine-MIL Project}
\date{May 2026}

\begin{document}
\maketitle
\tableofcontents
\newpage

% ─────────────────────────────────────────────────────────────────────────────
\section{Background and Motivation}
% ─────────────────────────────────────────────────────────────────────────────

\subsection{Setting}

An AB-MIL model was trained to classify which cytokine was applied to a
pseudo-tube of PBMCs (91 classes: 90 cytokines + PBS). Each pseudo-tube
contains $\approx$480 cells stratified by 18 cell types across 12 donors
(Donor2 and Donor3 held out for validation).

Training proceeds in three stages:
\begin{enumerate}[noitemsep]
    \item \textbf{Stage 1}: Encoder pre-training with cell-type classification supervision.
    \item \textbf{Stage 2}: MIL fine-tuning with frozen encoder (20 epochs, LR $= 10^{-3}$).
    \item \textbf{Stage 3}: Full fine-tuning, encoder unfrozen (100 epochs).
\end{enumerate}
During Stage 3, per-cell-type centroids were logged at every epoch for all
cytokines and training donors. Experiments across \textbf{10 random seeds} were
run to assess robustness.

\subsection{Known Cascade Pairs (Benchmark)}

The PBS-RC static geometry method (preceding this session) recovers all 11 known
cascade pairs with high sensitivity, but calls \emph{every} pair as
``shared'' (both $A \to B$ and $B \to A$ significant). The open problem that
this session addresses is \textbf{directionality}: which cytokine is upstream?

\begin{table}[H]
\centering
\caption{Known cytokine cascade pairs used as benchmark throughout this session.}
\begin{tabular}{lll}
\toprule
\textbf{Source (A)} & \textbf{Target (B)} & \textbf{Known relay cell types} \\
\midrule
IL-12       & IFN-$\gamma$  & NK, NKT, CD8 T \\
IL-1-beta   & IL-6          & Monocytes, macrophages \\
IL-2        & IL-15         & T cells, NK \\
IL-33       & IL-13         & ILC2, mast cells \\
IL-18       & IFN-$\gamma$  & NK, T cells \\
IL-21       & IL-10         & B cells, plasma cells \\
TNF-alpha   & IL-6          & Stromal, monocytes \\
IFN-alpha1  & IFN-$\gamma$  & NK, T cells \\
IL-10       & IL-6          & Monocytes (regulatory) \\
IL-4        & IL-13         & Th2, ILC2 \\
IL-27       & IFN-$\gamma$  & NK, T cells \\
\bottomrule
\end{tabular}
\end{table}

\subsection{Model Performance}

The trained model achieves $P_{\text{correct}} \approx 0.039$ on training data
(3.5$\times$ above chance for 91 classes). This relatively low accuracy is
an important caveat for interpreting all experiments in this session.

% ─────────────────────────────────────────────────────────────────────────────
\section{Experiment 1: Attention-Weighted Centroid PCA Trajectories}
% ─────────────────────────────────────────────────────────────────────────────

\subsection{Hypothesis}

\textbf{If A$\to$B via relay cell type T}, the model should learn to represent
T cells in A-tubes as B-like. Under fine-tuning, the attention-weighted
centroid of T cells in A-tubes ($\tilde{\mu}_{A,T}^{\text{attn}}$) should
drift \emph{toward} the centroid of T cells in B-tubes in the latent space.
Projecting both cytokines' centroids into a shared PCA 2D subspace (fitted on
all centroids from source and target across all epochs) should reveal this
convergence as a trajectory plot.

\subsection{Method}

For each seed and each cell type $T$:
\begin{enumerate}[noitemsep]
    \item Compute attention-weighted centroid: $\tilde{\mu}_{A,T}(t) = \text{mean}_{i:\tau(i)=T}\, a_i(t)\cdot \tilde{h}_i(t)$, averaged over training donors.
    \item Fit PCA on all centroids of source and target (all cell types, all epochs) for that seed.
    \item Project source and target centroids into the 2D cascade-relevant subspace.
    \item Track the trajectory over training epochs.
\end{enumerate}

\textbf{Analysis:} 10 seeds $\times$ 18 cell types, focused on IL-12 $\to$ IFN-$\gamma$.

\subsection{Results and Failure Mode}

Per-seed visualizations showed inconsistent patterns across seeds. The key
failure mode was identified: \textbf{the attention-weighted centroid is
contaminated by attention magnitude}. For low-attention cell types (e.g., B
cells), $a_i \approx 0$ for all cells, so the weighted centroid collapses to
$\approx \mathbf{0}$ for \emph{both} cytokines, producing spurious convergence.
For high-attention relay cells, the model pushes their embeddings to be
\emph{maximally discriminative} (far from each other), not similar — producing
apparent divergence.

\begin{figure}[H]
    \centering
    \includegraphics[width=\textwidth]{figures/pca_distance_heatmap_IL_12_IFN_gamma_10seeds.png}
    \caption{%
        \textbf{PCA distance convergence heatmap} (IL-12 $\to$ IFN-$\gamma$,
        10 seeds $\times$ 18 cell types).
        Metric: $(d_{\text{start}} - d_{\text{end}}) / d_{\text{start}}$ in the
        per-seed cascade-relevant 2D PCA subspace; clipped to $\pm 3$.
        Green = converged (source and target centroids moved closer).
        Red = diverged. \textbf{Sign test (right panel):}
        B Naive and B Intermediate converge in $>$80\% of seeds --- biologically
        incorrect (B cells are not relay cells). NK, NKT converge in only
        $\approx$40--50\% of seeds (near chance), despite being the known relays.
        The metric is fundamentally unreliable.
    }
    \label{fig:pca_heatmap}
\end{figure}

\subsection{Conclusion}

\textbf{Null result.} The attention-weighted centroid convergence metric does not
reliably identify relay cells or cascade direction. The attention weighting
itself is the confound: low-attention cells produce trivial convergence; relay
cells (high attention) diverge because the model makes them discriminative.

% ─────────────────────────────────────────────────────────────────────────────
\section{Experiment 2: Epoch-Dependent PBS-RC Directional Bias Asymmetry}
% ─────────────────────────────────────────────────────────────────────────────

\subsection{Hypothesis}

\textbf{If A$\to$B via relay T}, the attention-weighted centroid of T cells in
A-tubes should deviate \emph{toward} B's global embedding, more than the
symmetric reverse signal. Define:
\begin{align}
b_{\text{fwd}}(A \to B, T, t) &= (\tilde{\mu}_{A,T}(t) - \tilde{\mu}_A(t)) \cdot \hat{u}_{A\to B}(t), \\
b_{\text{fwd}}(B \to A, T, t) &= (\tilde{\mu}_{B,T}(t) - \tilde{\mu}_B(t)) \cdot \hat{u}_{B\to A}(t),
\end{align}
where $\tilde{\mu}_A(t)$ is A's global centroid (mean over all cell types) and
$\hat{u}_{A\to B}(t)$ is the unit direction from A to B at epoch $t$.
The asymmetry $\text{Asym}(A\to B, T, t) = b_{\text{fwd}}(A\to B) - b_{\text{fwd}}(B\to A)$
should be \textbf{positive and growing} for cascade pairs.

\subsection{Method}

Computed from existing \texttt{dynamics\_stage3.pkl} files (no new training).
For each of the 11 known pairs, computed $b_{\text{fwd}}$ at every logged
epoch across all training donors and 10 seeds. Summarised by mean asymmetry
over epochs and cell types.

\subsection{Results}

\begin{table}[H]
\centering
\caption{%
    Epoch asymmetry summary. Mean Asym $= \overline{\text{Asym}(A\to B, T, t)}$
    averaged over all cell types and epochs. Call: A$\to$B if Mean Asym $> 0$.
    $\checkmark$ = matches known direction; $\times$ = does not match.
}
\begin{tabular}{llrll c}
\toprule
\textbf{A} & \textbf{B} & \textbf{Mean Asym} & \textbf{Best relay T} & \textbf{Call} & \\
\midrule
IL-12      & IFN-$\gamma$ & $-0.0028$ & MAIT       & B$\to$A & \textcolor{wrong}{$\times$} \\
IL-2       & IL-15        & $+0.0443$ & cDC        & A$\to$B & \textcolor{correct}{$\checkmark$} \\
IL-33      & IL-13        & $+0.0102$ & cDC        & A$\to$B & \textcolor{correct}{$\checkmark$} \\
IL-18      & IFN-$\gamma$ & $+0.0537$ & CD14 Mono  & A$\to$B & \textcolor{correct}{$\checkmark$} \\
IL-21      & IL-10        & $-0.0040$ & NKT        & B$\to$A & \textcolor{marginal}{?} \\
IL-10      & IL-6         & $-0.0210$ & ILC        & B$\to$A & \textcolor{wrong}{$\times$} \\
IL-4       & IL-13        & $+0.0223$ & Plasmablast & A$\to$B & \textcolor{correct}{$\checkmark$} \\
IL-27      & IFN-$\gamma$ & $-0.0064$ & CD14 Mono  & B$\to$A & \textcolor{wrong}{$\times$} \\
\bottomrule
\multicolumn{6}{l}{\small 3 pairs absent from trajectory data (IL-1-beta, TNF-alpha, IFN-alpha1).} \\
\multicolumn{6}{l}{\small \textbf{Score: 4/8 correct (50\% = chance).}} \\
\end{tabular}
\label{tab:epoch_asym}
\end{table}

\begin{figure}[H]
    \centering
    \begin{subfigure}[t]{0.48\textwidth}
        \includegraphics[width=\textwidth]{figures/epoch_asym_IL_12_IFN_gamma.png}
        \caption{IL-12 $\to$ IFN-$\gamma$: essentially flat asymmetry across all
        100 epochs. MAIT identified as best relay (biologically incorrect; NK
        cells are the known relay). Mean Asym $= -0.003 \approx 0$.}
    \end{subfigure}
    \hfill
    \begin{subfigure}[t]{0.48\textwidth}
        \includegraphics[width=\textwidth]{figures/epoch_asym_IL_18_IFN_gamma.png}
        \caption{IL-18 $\to$ IFN-$\gamma$: positive and growing asymmetry
        (CD14 Mono, CD4 Naive). Mean Asym $= +0.054$. This is what a genuine
        directional signal would look like, but the relay cell type is
        biologically unexpected.}
    \end{subfigure}
    \caption{%
        Epoch-dependent directional bias asymmetry for two representative pairs.
        \textbf{Left panel of each figure}: $b_{\text{fwd}}$ for A-tubes (solid)
        and B-tubes (dashed) per cell type over training epochs.
        \textbf{Right panel}: $\text{Asym}(A\to B, T, t)$; positive = A$\to$B signal.
        Mean across 10 seeds.
    }
    \label{fig:epoch_asym}
\end{figure}

\subsection{Conclusion}

\textbf{Null result} for cascade directionality (4/8 = chance). The canonical
IL-12$\to$IFN-$\gamma$ pair shows zero signal. Asymmetry magnitudes are
uniformly small ($|$Asym$| < 0.06$). The same root cause as Experiment 1
applies: the attention-weighted centroid confounds geometry with attention
magnitude.

% ─────────────────────────────────────────────────────────────────────────────
\section{Experiment 3: Confusion Matrix Asymmetry}
% ─────────────────────────────────────────────────────────────────────────────

\subsection{Hypothesis}

\textbf{If A$\to$B via relay cells}, A-tubes contain cells that transcriptionally
resemble B (they are producing B in response to A). The trained 91-class
classifier should therefore assign higher softmax probability to class B when
evaluating A-tubes than it assigns to class A when evaluating B-tubes:
\[
P(B \mid A\text{-tube}) > P(A \mid B\text{-tube}) \implies A \to B.
\]
This defines the asymmetry:
$\text{Asym}(A \to B) = C[A,B] - C[B,A]$, where $C[i,j] = \overline{P(j \mid \text{tube of class } i)}$.

\subsection{Method}

For each of the 10 seeds, loaded \texttt{model\_stage3.pt}, ran a forward pass
on all 9{,}100 training tubes, and accumulated the mean-softmax 91$\times$91
confusion matrix. Averaged across seeds and computed element-wise asymmetry.

\subsection{Results}

\begin{figure}[H]
    \centering
    \includegraphics[width=0.82\textwidth]{figures/confusion_asymmetry.png}
    \caption{%
        \textbf{Full 90$\times$90 confusion asymmetry matrix} (PBS excluded).
        $\text{Asym}[A,B] = P(B \mid A\text{-tube}) - P(A \mid B\text{-tube})$;
        blue = A$\to$B signal. The matrix is dominated by \textbf{random noise}
        --- no block structure or interpretable pattern is visible. The colorbar
        range is $\approx \pm 0.015$, consistent with fluctuations around zero.
    }
    \label{fig:conf_asym}
\end{figure}

\begin{table}[H]
\centering
\caption{Confusion asymmetry for known cascade pairs. All values are $< 0.005$ in magnitude (noise floor $\approx 1/91 \approx 0.011$). Score: \textbf{4/11 = 36\%} (below chance).}
\begin{tabular}{llrrl c}
\toprule
\textbf{A} & \textbf{B} & \textbf{Asym(A$\to$B)} & \textbf{Asym(B$\to$A)} & \textbf{Call} & \\
\midrule
IL-12      & IFN-$\gamma$ & $-0.00035$ & $+0.00035$ & B$\to$A & \textcolor{wrong}{$\times$} \\
IL-1-beta  & IL-6         & $-0.00000$ & $+0.00000$ & B$\to$A & \textcolor{wrong}{$\times$} \\
IL-2       & IL-15        & $+0.00431$ & $-0.00431$ & A$\to$B & \textcolor{correct}{$\checkmark$} \\
IL-33      & IL-13        & $-0.00006$ & $+0.00006$ & B$\to$A & \textcolor{wrong}{$\times$} \\
IL-18      & IFN-$\gamma$ & $-0.00011$ & $+0.00011$ & B$\to$A & \textcolor{wrong}{$\times$} \\
IL-21      & IL-10        & $-0.00037$ & $+0.00037$ & B$\to$A & \textcolor{marginal}{?} \\
TNF-alpha  & IL-6         & $+0.00147$ & $-0.00147$ & A$\to$B & \textcolor{correct}{$\checkmark$} \\
IFN-alpha1 & IFN-$\gamma$ & $-0.00021$ & $+0.00021$ & B$\to$A & \textcolor{wrong}{$\times$} \\
IL-10      & IL-6         & $+0.00331$ & $-0.00331$ & A$\to$B & \textcolor{correct}{$\checkmark$} \\
IL-4       & IL-13        & $+0.00204$ & $-0.00204$ & A$\to$B & \textcolor{correct}{$\checkmark$} \\
IL-27      & IFN-$\gamma$ & $-0.00078$ & $+0.00078$ & B$\to$A & \textcolor{wrong}{$\times$} \\
\bottomrule
\end{tabular}
\end{table}

\subsection{Conclusion}

\textbf{Null result.} The confusion matrix is pure noise. None of the known
cascade pairs appear in the top-20 most asymmetric pairs globally (those are
all dominated by coincidental cytokine similarities unrelated to cascades).

The root cause: with $P_{\text{correct}} \approx 0.039$, the off-diagonal
confusion elements are each $\approx 1\%$. The asymmetries we seek are on the
order of $0.001$ (0.1 percentage points), well below the signal-to-noise floor.

% ─────────────────────────────────────────────────────────────────────────────
\section{Experiment 4: Cell-Type Ablation}
% ─────────────────────────────────────────────────────────────────────────────

\subsection{Hypothesis}

\textbf{If T is a relay for A$\to$B}, removing T cells from A-tubes should
reduce $P(B \mid A\text{-tube})$ — because T carries the B-like signal that
makes A-tubes look like B. The relay score quantifies this:
\[
\text{relay\_score}(A \to B, T) = P(B \mid \text{full tube}) - P(B \mid \text{tube} \setminus T).
\]
\textbf{Directional test:}
$\text{relay\_score}(A \to B, T^*) > \text{relay\_score}(B \to A, T^*) \implies A \to B.$

This test operates purely on model outputs and makes no assumptions about the
embedding geometry. Cell types are identified from \texttt{obs['cell\_type']}
in the original H5AD files.

\subsection{Method}

For each of the 10 seeds and each training-donor tube of cytokine $A$:
(1) run the model on the full tube; (2) for each cell type $T$, remove all
cells of type $T$ and re-run the model; (3) compute relay score.
All 11 known cascade pairs were evaluated in both directions.
Directionality call: the direction with the higher relay score for the best
relay cell type wins.

\subsection{Results}

\begin{figure}[H]
    \centering
    \includegraphics[width=\textwidth]{figures/ablation_known_pairs.png}
    \caption{%
        \textbf{Cell-type ablation relay scores} for all 11 known cascade pairs.
        Red bars: $\text{relay\_score}(A \to B, T)$; blue bars:
        $\text{relay\_score}(B \to A, T)$. Top 8 cell types per pair by
        forward score. A$\to$B is indicated when red $>$ blue for the top cell type.
    }
    \label{fig:ablation}
\end{figure}

\begin{table}[H]
\centering
\caption{%
    Cell-type ablation directionality calls. ``Fwd'' = relay\_score$(A\to B)$;
    ``Rev'' = relay\_score$(B \to A)$. Ratio $= $ Fwd/Rev.
    Only pairs with ratio $>2$ are considered reliable.
}
\begin{tabular}{llllrrl c}
\toprule
\textbf{A} & \textbf{B} & \textbf{Best relay T} & \textbf{Fwd} & \textbf{Rev} & \textbf{Ratio} & \textbf{Call} & \\
\midrule
IL-2       & IL-15        & cDC         & 0.00755 & 0.00121 & 6.2$\times$ & A$\to$B & \textcolor{correct}{$\checkmark$} \\
IL-12      & IFN-$\gamma$ & CD16 Mono   & 0.00119 & 0.00034 & 3.5$\times$ & A$\to$B & \textcolor{correct}{$\checkmark$} \\
TNF-alpha  & IL-6         & CD4 Naive   & 0.00630 & 0.00301 & 2.1$\times$ & A$\to$B & \textcolor{correct}{$\checkmark$} \\
IL-27      & IFN-$\gamma$ & CD14 Mono   & 0.00303 & 0.00192 & 1.6$\times$ & A$\to$B & \textcolor{correct}{$\checkmark$} \\
IFN-alpha1 & IFN-$\gamma$ & CD16 Mono   & 0.00242 & 0.00207 & 1.2$\times$ & A$\to$B & \textcolor{marginal}{$\checkmark$\textsuperscript{m}} \\
IL-4       & IL-13        & cDC         & 0.00546 & 0.00441 & 1.2$\times$ & A$\to$B & \textcolor{marginal}{$\checkmark$\textsuperscript{m}} \\
IL-33      & IL-13        & cDC         & 0.00012 & 0.00011 & 1.1$\times$ & A$\to$B & \textcolor{marginal}{$\checkmark$\textsuperscript{m}} \\
IL-10      & IL-6         & CD14 Mono   & 0.00734 & 0.00734 & 1.0$\times$ & B$\to$A & \textcolor{wrong}{$\times$} \\
IL-18      & IFN-$\gamma$ & CD16 Mono   & 0.00007 & 0.00020 & 0.4$\times$ & B$\to$A & \textcolor{wrong}{$\times$} \\
IL-21      & IL-10        & NK CD56br.  & 0.00003 & 0.00027 & 0.1$\times$ & B$\to$A & \textcolor{marginal}{?} \\
IL-1-beta  & IL-6         & MAIT        & 0.00000 & 0.00000 & ---         & B$\to$A & \textcolor{wrong}{$\times$} \\
\bottomrule
\multicolumn{8}{l}{\small \textsuperscript{m} Marginal (ratio $<1.5\times$, within noise). 7/11 A$\to$B calls; 3/3 high-confidence ($>2\times$) correct.}
\end{tabular}
\end{table}

\subsection{Conclusion}

\textbf{Partial signal.} Of 3 high-confidence calls (ratio $>2\times$), all
3 are correct: IL-2$\to$IL-15 (6.2$\times$), IL-12$\to$IFN-$\gamma$ (3.5$\times$),
TNF-alpha$\to$IL-6 (2.1$\times$). However, scores are universally small
($< 0.01$), the relay cell types are not always biologically canonical (CD16
Mono rather than NK for IL-12), and 4 pairs show no or reversed signal.
The same low-accuracy bottleneck ($P_{\text{correct}} \approx 4\%$) limits
the signal-to-noise ratio throughout.

% ─────────────────────────────────────────────────────────────────────────────
\section{Cross-Experiment Comparison}
% ─────────────────────────────────────────────────────────────────────────────

\begin{table}[H]
\centering
\caption{%
    Summary of directional calls across all three experiments for each
    known cascade pair. \textcolor{correct}{$\checkmark$} = correct;
    \textcolor{wrong}{$\times$} = incorrect; \textcolor{marginal}{$\sim$} = marginal/uncertain.
    Bolded pairs are those where at least two experiments agree on the correct direction.
}
\begin{tabular}{ll ccc}
\toprule
\textbf{A} & \textbf{B} &
\textbf{Exp.\ 2 (Epoch Asym.)} &
\textbf{Exp.\ 3 (Conf.\ Matrix)} &
\textbf{Exp.\ 4 (Ablation)} \\
\midrule
\textbf{IL-2}      & \textbf{IL-15}       & \textcolor{correct}{$\checkmark$} & \textcolor{correct}{$\checkmark$} & \textcolor{correct}{$\checkmark$} \\
\textbf{IL-4}      & \textbf{IL-13}       & \textcolor{correct}{$\checkmark$} & \textcolor{correct}{$\checkmark$} & \textcolor{marginal}{$\sim$} \\
\textbf{TNF-alpha} & \textbf{IL-6}        & ---       & \textcolor{correct}{$\checkmark$} & \textcolor{correct}{$\checkmark$} \\
IL-33      & IL-13        & \textcolor{correct}{$\checkmark$} & \textcolor{wrong}{$\times$}  & \textcolor{marginal}{$\sim$} \\
IL-18      & IFN-$\gamma$ & \textcolor{correct}{$\checkmark$} & \textcolor{wrong}{$\times$}  & \textcolor{wrong}{$\times$} \\
IL-12      & IFN-$\gamma$ & \textcolor{wrong}{$\times$}  & \textcolor{wrong}{$\times$}  & \textcolor{correct}{$\checkmark$} \\
IL-27      & IFN-$\gamma$ & \textcolor{wrong}{$\times$}  & \textcolor{wrong}{$\times$}  & \textcolor{correct}{$\checkmark$} \\
IFN-alpha1 & IFN-$\gamma$ & ---       & \textcolor{wrong}{$\times$}  & \textcolor{marginal}{$\sim$} \\
IL-10      & IL-6         & \textcolor{wrong}{$\times$}  & \textcolor{correct}{$\checkmark$} & \textcolor{wrong}{$\times$} \\
IL-21      & IL-10        & \textcolor{marginal}{?}       & \textcolor{marginal}{?}       & \textcolor{marginal}{?} \\
IL-1-beta  & IL-6         & ---       & \textcolor{wrong}{$\times$}  & \textcolor{wrong}{$\times$} \\
\midrule
\textbf{Score} & & 4/8 (50\%) & 4/11 (36\%) & 7/11 (64\%) \\
\bottomrule
\end{tabular}
\end{table}

% ─────────────────────────────────────────────────────────────────────────────
\section{Overall Conclusions}
% ─────────────────────────────────────────────────────────────────────────────

\subsection{What Did Not Work}

\begin{enumerate}
    \item \textbf{Attention-weighted centroid convergence / PCA distance} (Exp.\ 1):
    Fundamentally flawed metric. The attention weighting conflates embedding
    geometry with the model's confidence routing. Relay cells are pushed to
    \emph{discriminative} positions, not B-like positions.

    \item \textbf{Epoch-dependent directional bias} (Exp.\ 2): Inherits the
    same flaw. The signal is near-zero for most pairs; 4/8 = chance.

    \item \textbf{Confusion matrix asymmetry} (Exp.\ 3): The 91-class model's
    low accuracy makes off-diagonal elements indistinguishable from noise.
    The known cascade pairs do not rank among the top-20 globally asymmetric pairs.
\end{enumerate}

\subsection{What Showed Partial Promise}

\textbf{Cell-type ablation} (Exp.\ 4) is the most principled approach:
it makes no assumption about embedding geometry and operates directly on model
outputs. The 3 high-confidence calls (ratio $>2\times$) are all biologically
correct. However, the signal strength is limited by model accuracy.

\subsection{Root Cause: Model Accuracy}

All four experiments share a bottleneck: $P_{\text{correct}} \approx 0.039$
on a 91-class problem. At this accuracy level:
\begin{itemize}[noitemsep]
    \item Off-diagonal confusion elements are each $\approx 1/91 \approx 1\%$
    \item Asymmetries of $\approx 0.001$ (0.1 pp) are required to detect cascade direction
    \item This is well within the noise floor for both the confusion matrix and ablation methods
\end{itemize}

\subsection{Path Forward}

The conceptually cleanest next step is to revisit the \textbf{PBS-RC static
geometry} (which already recovers 11/11 pairs with Wilcoxon significance) and
address directionality by comparing effect sizes rather than binary significance
calls. Alternatively, the \textbf{cell-type ablation} approach with a
higher-accuracy model (e.g., more training epochs or a smaller cytokine subset)
could provide reliable directional calls.

\vspace{2em}
\hrule
\vspace{0.5em}
\small
\textit{All experiments used 10 random seeds (seed values: 42, 123, 7, 17, 31,
53, 97, 137, 191, 257), Stage 3 model checkpoint (100 epochs full fine-tuning),
attention-weighted centroid trajectory, training donors only (Donor2, Donor3
held out). Code: \texttt{github.com/Yam-Arieli/cytokine-mil}.}

\end{document}
